\section*{Chapitre \ref{ch:b2:mantle-melting} --- Fusion du manteau}

\begin{solution}{exo:b2:mantle-melting:1}
(a) Plus de \qty{90}{\percent} d'olivine~: \omterm{def:b2:mantle-melting:peridotites}{dunite}. (b) De l'olivine et de l'orthopyroxène
sans clinopyroxène~: \omterm{def:b2:mantle-melting:peridotites}{harzburgite}. (c) Plus de \qty{40}{\percent} d'olivine avec
les deux pyroxènes~: \omterm{def:b2:mantle-melting:peridotites}{lherzolite}. (d) Moins de \qty{40}{\percent} d'olivine~: une pyroxénite
(une olivine-websterite, avec les deux pyroxènes). (e) $57/98 = \qty{58.2}{\percent}$,
$28/98 = \qty{28.6}{\percent}$, $13/98 = \qty{13.3}{\percent}$~: une \omterm{def:b2:mantle-melting:peridotites}{lherzolite}.
\end{solution}

\begin{solution}{exo:b2:mantle-melting:2}
(a) Le liquidus à \qty{60}{\percent} d'anorthite est à \qty{1390}{\celsius}, dans le domaine
de l'anorthite~: l'anorthite cristallise en premier. (b) Le diopside apparaît à
l'eutectique, \qty{1270}{\celsius}, lorsque le liquide a \qty{42}{\percent} d'anorthite.
(c) Le premier liquide de tout mélange des deux minéraux est le liquide eutectique,
à \qty{42}{\percent} d'anorthite, à \qty{1270}{\celsius}.
\end{solution}

\begin{solution}{exo:b2:mantle-melting:3}
(a) $\num{4.0e-5}\times1623/(3300\times1000) = \qty{1.97e-8}{K/Pa} =
\qty{19.7}{K/GPa}$. (b) $\num{e9}/(3300\times9.81) = \qty{3.09e4}{m}$, soit
\qty{30.9}{km} par gigapascal~; $19.7/30.9 = \qty{0.64}{K/km}$, environ cinquante fois
moins que les \qty{30}{K/km} près de la surface, où la chaleur est conduite et non
transportée.
\end{solution}

\begin{solution}{exo:b2:mantle-melting:4}
(a) $T_{\mathrm s}(0) = \qty{1085.7}{\celsius}$~; $T_{\mathrm s}(1) = 1085.7 +
132.9 - 5.1 = \qty{1213.5}{\celsius}$~; $T_{\mathrm s}(2) = 1085.7 + 265.8 - 20.4 =
\qty{1331.1}{\celsius}$~; $T_{\mathrm s}(3) = 1085.7 + 398.7 - 45.9 =
\qty{1438.5}{\celsius}$. (b) $132.9 - 2\times5.1\times2 = \qty{112.5}{K/GPa}$.
(c) À \qty{3}{GPa}, \qty{1400}{\celsius} est sous le solidus,
\qty{1438.5}{\celsius}~; en remontant à \qty{2}{GPa} le long de l'adiabate, elle se refroidit de
seulement environ \qty{20}{\celsius}, jusqu'à \qty{1380}{\celsius} environ, tandis que le solidus descend à
\qty{1331.1}{\celsius}~: la roche est maintenant au-dessus de son solidus et une partie fond.
\end{solution}

\begin{solution}{exo:b2:mantle-melting:5}
(a) $1350 + (1277.5 - 1350)\times(30 - 20)/(40 - 20) = \qty{1314}{\celsius}$.
(b) La composition se situe du côté diopside de l'eutectique~: diopside.
(c) Le liquide atteint l'eutectique avec \qty{42}{\percent} d'anorthite~; l'anorthite
reste dans le liquide, donc la fraction liquide vaut $30/42 = 0.714$ et $1 - 0.714 =
\qty{29}{\percent}$ du système a cristallisé en diopside. (d) Tout cristallise~:
\qty{70}{\percent} de diopside et \qty{30}{\percent} d'anorthite, la
composition globale.
\end{solution}

\begin{solution}{exo:b2:mantle-melting:6}
(a) Forstérite~: $60.08/(2\times40.30 + 60.08) = 60.08/140.68 =
\qty{42.71}{\percent}$ \ce{SiO2}~; enstatite~: $60.08/100.38 = \qty{59.85}{\percent}$.
(b) $0.975\times59.85 + 2.5 = \qty{60.85}{\percent}$. (c) L'enstatite se situe entre
la forstérite et le liquide~: la fraction de forstérite vaut $(60.85 - 59.85)/(60.85 -
42.71) = 1.00/18.14 = \qty{5.5}{\percent}$, valeur de Bowen et Andersen. (d) Le
liquide est plus riche en silice que l'orthopyroxène dont il provient, et le résidu
gagne de l'olivine.
\end{solution}

\begin{solution}{exo:b2:mantle-melting:7}
(a) $5.1P^2 - 112.9P + 214.3 = 0$~: $P_0 = (112.9 - \sqrt{12746.4 -
4371.7})/10.2 = (112.9 - 91.51)/10.2 = \qty{2.10}{GPa}$, soit $2.097\times30.9 =
\qty{64.8}{km}$. (b) $5.1P^2 - 112.9P + 364.3 = 0$~: $P_0 = (112.9 -
72.90)/10.2 = \qty{3.92}{GPa}$, soit \qty{121}{km}. (c) $T_{\mathrm s}'(2.10) = 132.9
- 21.4 = \qty{111.5}{K/GPa}$, donc $\Delta P_0 = 10/(111.5 - 20) = \qty{0.109}{GPa}$
par \qty{10}{\celsius}, soit \qty{3.4}{km}.
\end{solution}

\begin{solution}{exo:b2:mantle-melting:8}
(a) $T_{\mathrm s} = \qty{1331.1}{\celsius}$, $T_{\mathrm l} = 1475.0 + 160.0 -
12.8 = \qty{1622.2}{\celsius}$~; $T = 1331.1 + 0.10\times291.1 =
\qty{1360.2}{\celsius}$, soit \qty{1633}{K}. (b) $(\partial T/\partial P)_F = 112.5
+ 0.10\times(67.2 - 112.5) = \qty{108.0}{K/GPa}$~; $(\partial T/\partial F)_P =
T_{\mathrm l} - T_{\mathrm s} = \qty{291.1}{K}$. (c) $c_p/T = 0.612$~; avec $\alpha/\rho
= \qty{12.1}{J/(kg.K)}$ par gigapascal, le numérateur vaut $0.612\times108.0 - 12.1
= 54.0$ et le dénominateur $300 + 0.612\times291.1 = 478.2$~: $-\mathrm dF/\mathrm
dP = 0.113$, \qty{11.3}{\percent} par gigapascal. (d) $54.0/178.2 = 0.303$,
\qty{30}{\percent} par gigapascal~: la chaleur latente divise la productivité par
un facteur 2.7.
\end{solution}

\begin{solution}{exo:b2:mantle-melting:9}
(a) $43\times1 = \qty{43}{\celsius}$~; $43\times5^{0.75} = 43\times3.34 =
\qty{144}{\celsius}$~; $43\times13^{0.75} = 43\times6.85 = \qty{294}{\celsius}$.
(b) $1213.5 - 294 = \qty{919}{\celsius}$. (c) \qty{125}{ppm} correspondent à
\qty{0.0125}{\percent}, donc $X = 0.0125/0.01 = \qty{1.25}{\percent}$ et $\Delta T =
43\times1.25^{0.75} = 43\times1.18 = \qty{51}{\celsius}$.
\end{solution}

\begin{solution}{exo:b2:mantle-melting:10}
(a) La base a une largeur de $2H$ et le manteau la traverse à $U$~: $2UH$ par unité de
longueur de dorsale et de temps. Chaque côté couvre la hauteur $H$ et est traversé horizontalement à
$U$~: $UH$ chacun, $2UH$ au total. Avec des côtés à $\theta$, la base serait de
largeur $2H/\tan\theta$ et l'apport $2UH/\tan\theta$, égal à la sortie seulement si
$\tan\theta = 1$~: dans ce modèle, l'angle est fixé par la conservation de la masse. (b) Avec une
productivité constante, le \omterm{def:b2:mantle-melting:fraction}{taux de fusion} au sommet vaut $F_{\max} = \Pi P_0 = \Pi\rho
g z_0$, et $H = z_0$~: $h_c = F_{\max}z_0/2 = \Pi\rho g z_0^2/2$. (c) Chaque année,
\qty{50}{mm} de nouvelle plaque, $\qty{0.05}{m}\times\qty{1000}{m}\times\qty{7100}{m} =
\qty{3.6e5}{m^3}$ de croûte par kilomètre de dorsale~; avec un \omterm{def:b2:mantle-melting:fraction}{taux de fusion} moyen de
\qty{12}{\percent}, elle provient de $\num{3.6e5}/0.12 = \qty{3.0e6}{m^3}$ de manteau.
\end{solution}

\begin{solution}{exo:b2:mantle-melting:11}
(a) Après qu'une masse $t$ de A a quitté une masse unitaire de liquide, le liquide contient $0.6 -
t$ de A, 0.2 de B et 0.2 de C dans une masse $1 - t$~: $\big((0.6 - t), 0.2,
0.2\big)/(1 - t)$, et B/C reste égal à 1. (b) Sur la cotectique, $(0.40 - 0.15s, 0.60 -
0.30s, 0.45s)$~; B = C donne $0.60 - 0.30s = 0.45s$, $s = 0.8$~: $L_1$ contient A 28, B
36, C 36. Alors $0.2/(1 - t) = 0.36$, $t = 0.44$~: \qty{44}{\percent} du système a
cristallisé en A. (c) Le C est tout entier dans le liquide~: $f = 20/45 = 0.44$,
\qty{44}{\percent} encore liquide. (d) A~: $0.60 - 0.444\times0.25 = 0.49$~; B~: $0.20 -
0.444\times0.30 = 0.07$~; avec le liquide, $49 + 7 + 44 = \qty{100}{\percent}$.
\end{solution}

\begin{solution}{exo:b2:mantle-melting:12}
(a) $X_{\mathrm{sat}}(1) = 12 + 1 = \qty{13}{\percent}$, $\Delta T =
\qty{294}{\celsius}$, solidus saturé $1213.5 - 294 = \qty{919}{\celsius}$~;
$X_{\mathrm{sat}}(2) = 12\times1.516 + 2 = \qty{20.2}{\percent}$, $\Delta T =
43\times20.2^{0.75} = 43\times9.53 = \qty{410}{\celsius}$, solidus saturé $1331.1
- 410 = \qty{921}{\celsius}$. (b) Le liquide dissout davantage d'eau à plus haute pression,
et l'abaissement plus grand annule la hausse de \qty{118}{\celsius} du solidus sec. (c) Le solidus doit
baisser de $1331.1 - 1100 = \qty{231}{\celsius}$~: $X =
(231/43)^{1/0.75} = 5.37^{1.333} = \qty{9.4}{\percent}$ d'eau, inférieur aux
\qty{20.2}{\percent} à saturation, si bien que le coin peut fondre.
\end{solution}

\begin{solution}{pb:b2:mantle-melting:1}
\textbf{1.} $1350 + 273 = \qty{1623}{K}$.
\textbf{2.} $\num{4.0e-5}\times1623/(3300\times1000) = \qty{1.97e-8}{K/Pa} =
\qty{19.7}{K/GPa}$, arrondi à \qty{20}{K/GPa}.
\textbf{3.} $\num{e9}/(3300\times9.81) = \qty{30.9}{km}$~; $20/30.9 =
\qty{0.65}{K/km}$.
\textbf{4.} $1350 + 20P = 1085.7 + 132.9P - 5.1P^2$, soit $5.1P^2 - 112.9P +
264.3 = 0$: $P_0 = (112.9 - \sqrt{12746.4 - 5391.7})/10.2 = (112.9 - 85.76)/10.2 =
\qty{2.66}{GPa}$.
\textbf{5.} $2.66\times30.9 = \qty{82}{km}$~; $T = 1350 + 20\times2.66 =
\qty{1403}{\celsius}$.
\textbf{6.} $1085.7 + 132.9\times2.66 - 5.1\times2.66^2 = 1085.7 + 353.5 - 36.1 =
\qty{1403}{\celsius}$.
\textbf{7.} $T_{\mathrm s} = 1085.7 + 199.4 - 11.5 = \qty{1273.6}{\celsius}$~;
$T_{\mathrm l} = 1475.0 + 120.0 - 7.2 = \qty{1587.8}{\celsius}$~; différence
\qty{314.2}{K}.
\textbf{8.} $1273.6 + 0.10\times314.2 = \qty{1305.0}{\celsius}$, \qty{1578}{K}.
\textbf{9.} $T_{\mathrm s}' = 132.9 - 15.3 = \qty{117.6}{K/GPa}$, $T_{\mathrm l}' =
80.0 - 9.6 = \qty{70.4}{K/GPa}$~: $(\partial T/\partial P)_F = 117.6 +
0.10\times(70.4 - 117.6) = \qty{112.9}{K/GPa}$~; $(\partial T/\partial F)_P =
\qty{314.2}{K}$.
\textbf{10.} $c_p/T = 1000/1578 = 0.634$~; numérateur $0.634\times112.9 - 12.1 = 71.6
- 12.1 = 59.5$~; dénominateur $300 + 0.634\times314.2 = 499.2$~: $0.119$, soit
\qty{11.9}{\percent} par gigapascal.
\textbf{11.} Avec la fonction ajustée $F = T'^{1.5}$, où $T'$ est la fraction du chemin parcouru du solidus au liquidus, la température doit monter fortement avec $F$ pour produire
le premier pour cent de liquide, donc $(\partial T/\partial F)_P$ est grand près du solidus
et la fusion démarre lentement. $F_{\max} = 0.09\times2.66 = 0.24$, \qty{24}{\percent}.
\textbf{12.} $0.13/0.50 = 0.26$~: le clinopyroxène est épuisé à \qty{26}{\percent} de
fusion. À \qty{24}{\percent}, le résidu sous l'axe en garde encore un peu~:
une \omterm{def:b2:mantle-melting:peridotites}{lherzolite} pauvre en clinopyroxène, à la limite de la \omterm{def:b2:mantle-melting:peridotites}{harzburgite}.
\textbf{13.} $h_c = 0.24\times82/2 = \qty{9.8}{km}$.
\textbf{14.} $F_{\max}/2 = \qty{12}{\percent}$.
\textbf{15.} \qty{9.8}{km} dépassent de \qty{38}{\percent} les \qty{7.1}{km} mesurés, au-dessus de la
fourchette de la croûte océanique normale.
\textbf{16.} Une productivité constante donne au manteau le plus profond trop de liquide~; le
modèle ajusté fait apparaître le premier liquide lentement, $F(z)$ reste plus faible près du début de fusion, et
$\int F\,\mathrm dz$ est plus petite.
\textbf{17.} $F_{\max} = \Pi P_0$ et $H = 30.9P_0$ (km)~: $h_c = \Pi P_0\times30.9P_0/2
= kP_0^2$ avec $k = 0.09\times30.9/2 = \qty{1.39}{km/GPa^2}$.
\textbf{18.} $P_0 = \sqrt{7.1/1.39} = \sqrt{5.11} = \qty{2.26}{GPa}$.
\textbf{19.} $2.26\times30.9 = \qty{70}{km}$.
\textbf{20.} $T_{\mathrm s}(2.26) = 1085.7 + 300.4 - 26.0 = \qty{1360.0}{\celsius}$,
et $T_p = 1360.0 - 20\times2.26 = 1360.0 - 45.2 = \qty{1315}{\celsius}$.
\textbf{21.} Pour \qty{6.3}{km}~: $P_0 = \sqrt{4.53} = \qty{2.13}{GPa}$, $T_{\mathrm s}
= \qty{1345.6}{\celsius}$, $T_p = \qty{1303}{\celsius}$. Pour \qty{7.9}{km}~: $P_0 =
\sqrt{5.68} = \qty{2.38}{GPa}$, $T_{\mathrm s} = \qty{1373.1}{\celsius}$, $T_p =
\qty{1326}{\celsius}$. La dispersion de la croûte autorise $1315 \pm
\qty{12}{\celsius}$.
\textbf{22.} $P_0 = \sqrt{20/1.39} = \sqrt{14.39} = \qty{3.79}{GPa}$~; $T_{\mathrm s} =
1085.7 + 503.7 - 73.3 = \qty{1516.1}{\celsius}$~; $T_p = 1516.1 - 75.8 =
\qty{1440}{\celsius}$, soit un excès de $1440 - 1315 = \qty{125}{\celsius}$.
\textbf{23.} Les deux utilisent des indices différents. Les laves enregistrent le manteau le plus chaud,
au centre du panache~; la croûte moyenne les liquides de tout le manteau qui traverse la
région de fusion, dont une grande partie est plus froide. Le modèle simple garde aussi
la productivité constante, alors que le clinopyroxène est épuisé à \qty{26}{\percent}
tandis que $\Pi P_0 = 0.09\times3.79 = 0.34$ ici~: un manteau chaud réel fond moins, et il lui faut
une température plus élevée pour produire la même croûte.
\textbf{24.} $T_{\mathrm s}'(2.26) = 132.9 - 23.1 = \qty{109.8}{K/GPa}$~; $\Delta P_0 =
10/(109.8 - 20) = \qty{0.111}{GPa}$~; puisque $h_c \propto P_0^2$, $\Delta h_c/h_c =
2\times0.111/2.26 = \qty{9.8}{\percent}$, soit environ un dixième par \qty{10}{\celsius}.
\textbf{25.} Sous une dorsale très lente, le manteau monte si lentement qu'il perd de la chaleur par
conduction vers la surface froide~; il quitte l'adiabate, sa fusion s'arrête
plus profondément, et le sommet du triangle, où $F$ est le plus grand, est tronqué~: la croûte est
plus mince.
\textbf{26.} La croûte des dorsales normales donne une \omterm{def:b2:mantle-melting:tp}{température potentielle du manteau} de
\textbf{\qty[text-series-to-math=true]{1315}{\celsius}}, à environ
\qty{12}{\celsius} près~: \qty{35}{\celsius} sous les \qty{1350}{\celsius} de la physique des
minéraux, qui est incertaine d'environ autant, et bien à l'intérieur de l'intervalle de 1280 à
\qty{1400}{\celsius} issu des laves. Trois thermomètres indépendants concordent sur un manteau
voisin de \qty{1300}{\celsius} à \qty{1350}{\celsius}.
\end{solution}
